%=======================================================================
%  Paper V — Section 8 (Discussion)   FINAL
%
%  Requires the macros listed at the head of Section 2.
%  NOTE: ver10 used \dens in Eq. (macmillan). That macro is not defined
%  anywhere in this manuscript and would break the compile. It is
%  written out as \varrho below.
%
%  PRINCIPAL ADDITION IN THIS REVISION -------------------------------
%  Section 8.2 is new. It shows, from the mass constraint alone, that
%  the first condition of Richardson and Bowling -- that the rotational
%  term be a significant fraction of the gravitational one -- is the
%  condition under which THEIR objective is not identically flat, and
%  that it has no counterpart in a fixed-mass two-region inversion,
%  where the mass constraint fixes the scale instead. This converts the
%  discomfort acknowledged in Section 8.1 from an unresolved concern
%  into a bounded one, and it costs nothing: the omega = 0 control run
%  needs no new gravity evaluation.
%
%  Superseded values and statements:
%      "2e-14 per cent" volume closure  -> 2.1e-14 in relative terms
%      "2.2e-14 per cent" sphere limit  -> relative; and it is a check
%                                          on the analytic quadrature,
%                                          not on the polyhedral code
%      "19 to 59 times more accurate"   -> quoted at the working
%                                          resolution as well, ~16-20
%      "our 16.5 m lies within it"      -> 16.5 +/- 2.9 m, and the
%                                          Lowry interval contains the
%                                          whole sweep, so the test does
%                                          not discriminate
%      head-neck as the fundamental ambiguity
%                                       -> the method separates them by
%                                          a factor 7.3 in I; the real
%                                          degeneracy is elsewhere
%      "no numerical diagnostic detects a misplaced plane"
%                                       -> one of four fails outright,
%                                          two rank it correctly
%      "no change across the three Itokawa meshes"
%                                       -> bounded at half a scan step
%
%  Resolved in this pass: the three missing dispersion pairs of Table 8.2
%  (Itokawa does carry the largest residual, 83.1 % against 98.7 %); the
%  seven rungs of the planted ladder; the Paper IV slope range, 1.8-25.6 %
%  as published against 1.8-21.9 % recomputed with clustering here.
%  Also resolved: the omega = 0 control and the equatorial column
%  (v10_s8_tests.json), and the standard errors of Tables 8.3 and 8.4. The
%  49,152-facet row was abandoned as disproportionately expensive; see the
%  note at that table.
%=======================================================================

\section{Discussion}
\label{sec:discussion}

\subsection{The standing of the relaxation premise on these bodies}
\label{sec:discussion:premise}

Everything in this paper rests on a premise that is not established for any
of the bodies examined. \citet{richardson2014} state four conditions under
which a surface may be taken as relaxed toward an equipotential: the mean
rotational force must be a significant fraction of the mean gravitational
force; a sufficiently thick, low-cohesion, mobile regolith layer must cover
most of the surface; a downslope flow disturbance source must be present;
and sufficient time must have passed since the last major surface
alteration. They are explicit that the result is not a density measurement
but the density corresponding to the most eroded state of the body given its
present shape, topography and spin.

On Itokawa three of the four are supported and the first is not. Mobile
regolith is present: the Muses Sea and the other smooth terrains cover about
a fifth of the surface \citep{fujiwara2006,kanamaru2019pss}. A disturbance
source is present: \citet{tatsumi2018} invoke impact-induced seismic shaking
to explain the depletion of small quasi-circular depressions. Time is
available: they place the crater retention age at $3$--$33$~Myr in the main
belt, in agreement with the cosmic-ray exposure ages of the returned
samples, and conclude that global resurfacing occurred long after the
disruption of the parent body. The first condition, however, fails.
\citet{richardson2014} tabulate the mean rotational potential on Itokawa as
$2.4\%$ of the mean gravitational potential, the second lowest of the bodies
they examine, and their global minimization returns a bulk density of
$\SI{330}{\kilo\gram\per\cubic\metre}$ against a measured $1950$---an error
of a factor of six, which they attribute directly to that first criterion.
% Supporting evidence for "seven": the companion conference abstract
% (Richardson, Graves & Bowling, ACM 2014) states that the study used 22
% radar-derived shape models and SEVEN small-body shape models from
% spacecraft observations. The seven quoted here are the spacecraft set, the
% only ones with measured masses. Confirm against the Icarus paper itself
% before submission; the abstract is a companion, not the paper.

This is uncomfortable, and it should be stated rather than left for a reader
to find. The single strong signal reported here is on the body where the
premise is weakest, while Bennu, which rotates fast enough to satisfy the
condition comfortably, returns a weak-signal result.
Section~\ref{sec:discussion:conditioning} shows that the discomfort is
narrower than it first appears, because the condition that fails is a
condition on the conditioning of a different inverse problem; this
subsection first sets out where the premise stands body by body.

\begin{table}[htbp]
\centering
\small
\caption{The four conditions of \citet{richardson2014}, assessed body by
body. The rotation parameter in the second column is
$q = \omega^{2}R_{v}/(G\Mtot/R_{v}^{2})$, evaluated at the
volume-equivalent radius $R_{v}=(3\Vtot/4\pi)^{1/3}$, which reduces
identically to $q = 3\omega^{2}/(4\pi G\densbulk)$: it depends on the spin
period and the bulk density alone and carries no information about the
shape. It is therefore \emph{not} the ratio at the equator of an elongated
body, which is larger by the ratio of the maximum distance from the spin
axis to $R_{v}$---a factor of $1.89$ on Itokawa and $2.09$ on
Eros---and it is not the quantity \citet{richardson2014} tabulate, which is
a ratio of mean potentials. The three are listed or distinguished
separately and must not be conflated. Arrokoth has no measured mass, so its
$q$ is a function of the assumed bulk density and is quoted at all three
values run.}
\label{tab:conditions}
\setlength{\tabcolsep}{4.5pt}
\begin{tabular}{lcccccc}
\toprule
& \multicolumn{3}{c}{(1) rotation} & (2) mobile & (3) disturbance & (4) time \\
\cmidrule(lr){2-4}
Body & $q$ at $R_{v}$ & $q$ at equator & potential
 & regolith & source & since \\
\midrule
Itokawa  & $0.037$ & $0.070$ & $0.024$ & yes & yes & yes \\
Eros     & $0.147$ & $0.308$ & $0.111$ & yes & yes & yes \\
Bennu    & $0.496$ & $0.586$ & ---     & partly & yes & uncertain \\
67P      & $0.132$ & $0.211$ & ---     & partly & yes & no \\
Arrokoth & $0.173/0.108/0.086$ & $0.337/0.211/0.169$ & ---
         & unknown & weak & yes \\
\bottomrule
\end{tabular}

\vspace{4pt}
\begin{minipage}{\textwidth}
\footnotesize
The Arrokoth entries correspond to assumed bulk densities of $250$, $400$
and $\SI{500}{\kilo\gram\per\cubic\metre}$ respectively; at the central
published estimate near $\SI{235}{\kilo\gram\per\cubic\metre}$ the value is
$0.18$, which exceeds Eros's. The equatorial column is measured on the shape models used here, as
$\omega^{2}R_{\rm eq}/(G\Mtot/R_{v}^{2})$ with $R_{\rm eq}$ the largest
distance of any vertex from the spin axis: $306.1$ and $\SI{17.6}{\kilo\metre}$
on Itokawa and Eros, against $R_{v}$ of $161.9$ and
$\SI{8.43}{\kilo\metre}$, and ratios $R_{\rm eq}/R_{v}$ of $1.89$, $2.09$,
$1.18$, $1.60$ and $1.96$ for the five bodies in order.
\end{minipage}
% Measured on the meshes (v10_s8_tests.json): R_eq = 306.1 m, 17628.6 m,
% 289.4 m, 2633.5 m, 19482.7 m for Itokawa, Eros, Bennu, 67P and Arrokoth.
% The earlier Itokawa entry used a semi-axis of 267.5 m; the mesh gives 306.1.
\end{table}

The premise is not equally plausible across the five, and the reader should
be able to see where it is weak. Table~\ref{tab:conditions} sets the four
conditions against each body; the paragraphs below give the reasoning. One
caution applies throughout: \citet{richardson2014} do not state a numerical
threshold for ``a significant fraction'', so no body is judged against a
stated cut. The evidence that Itokawa falls below it is empirical---their own
bulk-density inversion on that body errs by a factor of six---and the
evidence that Eros lies above it is that the same inversion works there. The
threshold in mean potential ratio is therefore bracketed between $0.024$ and
$0.111$ and is not located.

\emph{Itokawa} satisfies three and fails the first, as described above.

\emph{Eros} is the most favourable case on the premise as a whole. Its
rotation parameter is four times Itokawa's, its regolith is deep enough to
have formed ponded deposits, impacts supply a disturbance source, and its
surface is old. That the method returns no contrast on the body where the
premise is strongest, and where the measured gravity field independently
indicates near-uniformity \citep{yeomans2000}, is the cleanest control in
the study.

\emph{Bennu} satisfies the first condition most comfortably of all five, at
a rotation parameter approaching one half. The other three are less clear.
Its surface is boulder-rich rather than smoothly regolith-covered, though
particles are observed to be ejected and to move, so material is mobile in
some sense. And its surface may be young: \citet{scheeres2020} attribute the
reported low-density equator to recent migration and redistribution of
material, which is a statement that the surface has changed recently rather
than settled. A surface still in motion is not a relaxed surface.

\emph{67P} fails the fourth condition outright. A comet nucleus is
resurfaced continuously by sublimation, and its shape is not the product of
gravitational relaxation alone. Whatever mobile material exists is airfall
dust redeposited between perihelion passages rather than a regolith that has
flowed downslope to equilibrium. This is a further reason, independent of
the model-class argument of \citet{jorda2016}, to expect no interpretable
signal.

\emph{Arrokoth} is the least constrained, and its entry in
Table~\ref{tab:conditions} illustrates why. Its rotation parameter is not a
measured quantity at all but a function of the assumed bulk density, running
from $0.173$ to $0.086$ across the range tried; an earlier draft described
it as comparable to Eros's, which is true only at the middle of that range,
and at the central published estimate it exceeds Eros's. Beyond that,
nothing is known of its regolith, no disturbance source is established for a
primordial Kuiper Belt object, and its surface is thought to have been
little modified since accretion \citep{mckinnon2020}---which satisfies the
fourth condition in the sense of elapsed time while removing any reason to
think relaxation ever operated. With no measured mass on top of this, its
result carries the least weight of the five.

The pattern is worth stating plainly. The premise is strongest on Eros,
which returns no signal; it is weakest on Itokawa, which returns the only
strong one. The four weak-signal cases are therefore not equally
informative: Eros's is a result about Eros, while 67P's and Arrokoth's are
as much about the premise as about the bodies.

Two of the three considerations that bear on how much weight this should
carry are empirical, and are taken up here; the third is structural and is
the subject of Section~\ref{sec:discussion:conditioning}.

The first empirical consideration is that the result does not stand alone,
and the second is that this support is weaker than it first appears.
\citet{lowry2014} derive a centre-of-mass offset of $\SI{21+-12}{\metre}$
from an observed spin-up, an argument that makes no appeal to surface
relaxation at all, and our $\SI{16.5+-2.9}{\metre}$ at the same dividing
plane lies within it. Two inferences resting on different premises agree,
and by Section~\ref{sec:constrains:dipole} they are genuinely independent
rather than two readings of one quantity, since the objective function
demonstrably does not hold the mass dipole fixed. But the agreement should
not be asked to carry more than it can. The Lowry determination is itself
only $1.75$ standard deviations from zero, so it does not exclude a
homogeneous Itokawa at the two-sigma level; and its interval,
$9$ to $\SI{33}{\metre}$, contains the centre-of-mass offset at
\emph{every} plane of Table~\ref{tab:itokawasweep}, whose range is
$14.89$ to $\SI{20.52}{\metre}$, together with the $\pm2.9$ interval on each.
The external check therefore confirms the sign and the order of magnitude of
the anomaly and discriminates neither the plane nor the value.

The third empirical consideration cuts the other way, and is recorded here.
Even after the best two-region fit, Itokawa's surface remains $7.2\%$ from
an equipotential in the normalized dispersion, against $3.4\%$ for Eros with
no contrast fitted at all. Three things should be said about that
comparison before it is read as evidence that Itokawa is the less relaxed.
It is not a resolution effect: a $5.9\%$ change in $\potvarn$, the top of
the Paper~IV envelope, is $0.2$ percentage points on Eros's residual against
the $3.8$ percentage points separating the two bodies. It is, however, not
controlled for shape, since a more irregular body has a larger $\potvarn$
whatever its state of relaxation, and Itokawa is a contact binary while Eros
is not. And it is incomplete: the residuals for 67P, Bennu and Arrokoth are
not reported anywhere in this paper, so the claim that Itokawa's is the
largest of the five is not yet supported. The missing entries are the
decisive ones, because 67P and Arrokoth are also bilobate: if their
residuals are comparable to Itokawa's, the $7.2\%$ is shape, and if they are
not, it is relaxation.

\begin{table}[htbp]
\centering
\small
\caption{Normalized geopotential dispersion before and after the two-region
fit, at the plane of Table~\ref{tab:planes} for each body. The final column
is the fraction of the departure from an equipotential that the best
two-region density model does \emph{not} remove, and it is large on every
body: on Itokawa, where the fit is by far the most effective, five-sixths of
the departure remains and is not of a form any model of this class can
absorb.}
\label{tab:residual}
\setlength{\tabcolsep}{10pt}
\begin{tabular}{lrrrr}
\toprule
Body & $\potvarn$ (uniform) & $\potvarn$ (best fit) & $\improve$
 & residual \\
\midrule
Itokawa  & $0.08663$ & $0.07200$ & $16.9\%$  & $83.1\%$ \\
Eros     & $0.03413$ & $0.03399$ & $0.418\%$ & $99.6\%$ \\
Bennu    & $0.06158$ & $0.06158$ & $<10^{-3}\%$ & $100.0\%$ \\
67P      & $0.07615$ & $0.07571$ & $0.574\%$ & $99.4\%$ \\
Arrokoth & $0.05679$ & $0.05607$ & $1.272\%$ & $98.7\%$ \\
\bottomrule
\end{tabular}
\end{table}
% Filled from paperV_final.json; the Arrokoth pair is from the run with the
% buried facets excluded (arrokoth_clean.json), matching the 1.272 % quoted
% throughout. Itokawa does carry the largest residual: 83.1 % against 98.7 %
% for the next body. The superseded note read:
% overturn it.

\subsection{Why the first condition does not transfer to this inversion}
\label{sec:discussion:conditioning}

The two inference problems are not the same, and the difference can be made
exact rather than asserted. \citet{richardson2014} vary the total mass,
which changes the balance between the gravitational and rotational terms;
this work holds the measured mass fixed and varies only the internal
distribution. That distinction is what the first condition turns on.

Consider their problem first. With a single uniform density,
Equation~\eqref{eq:geopotential} reads
$\Upot_i = \varrho\,\Ufull_i + \Urot_i$, so
%
\begin{equation}
\potvarn^{2}(\varrho)
 = \frac{\varrho^{2}\operatorname{Var}(\Ufull)
        + 2\varrho\operatorname{Cov}(\Ufull,\Urot)
        + \operatorname{Var}(\Urot)}
        {\left(\varrho\,\overline{\Ufull}+\overline{\Urot}\right)^{2}} ,
\label{eq:uniformobjective}
\end{equation}
%
where variances and covariances are area-weighted over the surface. If
$\Urot \equiv 0$ this collapses to
$\potvarn = \sqrt{\operatorname{Var}(\Ufull)}\big/\lvert\overline{\Ufull}\rvert$,
which is independent of $\varrho$: the objective is identically flat and the
bulk density is not determined at all. Their first condition is precisely
the condition that Equation~\eqref{eq:uniformobjective} is not flat, and a
rotational potential ratio of $0.024$, as on Itokawa, leaves a minimum too
shallow to locate---which is why their inversion there returns $330$ against
a measured $1950$.

Now substitute the mass constraint~\eqref{eq:massbalance} into
Equation~\eqref{eq:geopotential}. With
$a \equiv \Mtot/(\Vtot-\VR)$ and $b \equiv \VR/(\Vtot-\VR)$, so that
$\densrest = a - b\densR$, the surface geopotential is affine in the single
remaining free parameter,
%
\begin{equation}
\Upot_i = \alpha_i + \beta_i\,\densR ,
\qquad
\alpha_i = a\left(\Ufull_i - \UR_i\right) + \Urot_i ,
\qquad
\beta_i = (1+b)\,\UR_i - b\,\Ufull_i .
\label{eq:affine}
\end{equation}
%
Set $\Urot \equiv 0$, as the first condition's failure would in the limit
require. Then $\alpha_i = a(\Ufull_i - \UR_i)$, and $\alpha$ and $\beta$ are
proportional only if $a(1+b) = ab$, that is only if $a = 0$. Since
$a = \Mtot/(\Vtot-\VR)$ is strictly positive for any region smaller than the
body, they are never proportional, and $\potvarn(\densR)$ remains a
non-constant ratio of quadratics with an interior stationary point. The
degeneracy that the first condition protects against does not arise here,
because the mass constraint fixes the scale of the density field that the
rotational term fixes in their problem.

This is a statement about conditioning and not about physics, and the
distinction matters. It shows that a slow rotator does not make the
fixed-mass two-region problem ill-posed; it does not show that the surface of
a slow rotator has relaxed, which remains the premise and remains unverified.
What it removes is the specific inference that because Richardson and
Bowling's method fails on Itokawa by a factor of six, a method built on their
premise must fail there too. The two failures would have different causes,
and only one of the causes is present.

Two empirical checks bear on the argument and we report the first. A
fourfold change in the assumed total mass moves the recovered contrast on
Itokawa by $11\%$ and never near unity
(Section~\ref{sec:constrains:mass}), whereas the bulk-density inference on
the same body lands at $330$ against $1950$: the fixed-mass problem is
demonstrably better conditioned on this body than the variable-mass one, by
a wide margin. The second check is direct, and we have now run it. It confirms the
derivation in the form that matters and corrects the form in which an earlier
draft anticipated it. With $\Urot$ set identically to zero the scan returns
an interior minimum at
$\denshead = \SI{2950}{\kilo\gram\per\cubic\metre}$ and
$\densrest = \SI{1828.6}{\kilo\gram\per\cubic\metre}$, a contrast of
$1.613$ and an improvement of $17.16\%$, against $2800$, $1858.1$, $1.507$
and $16.89\%$ with the rotational term retained. The structural claim is
therefore established in the output and not only in the derivation: the
two-region objective does not collapse without rotation, whereas the
single-density objective of \citet{richardson2014} becomes identically flat
under the same substitution and determines nothing at all. What supplies the
scale there is rotation; what supplies it here is the mass constraint.

What the control does not show is that rotation is unimportant. The contrast
moves by $7.1\%$, which exceeds its own error budget of $6.0\%$, and the
background density by $\SI{29.5}{\kilo\gram\per\cubic\metre}$, which
exceeds the $\pm\SI{27}{\kilo\gram\per\cubic\metre}$ quoted for it in
Section~\ref{sec:itokawa}. A rotational potential that is $2.4\%$ of the
gravitational one still carries a measurable share of the inference, and it
is twice what the choice of spin axis contributes
(Section~\ref{sec:methods}). One detail of the control is worth recording:
the improvement is \emph{higher} without rotation, $17.16\%$ against
$16.89\%$. The centrifugal term contributes dispersion that no two-region
density model can remove, so its presence lowers the ceiling on what the fit
can reach rather than raising the floor. The correct summary is that a slow
rotator is not ill-conditioned for this inversion, not that rotation may be
neglected in it.
% Run (v10_s8_tests.json): omega = 0 gives rho_R 2950, rho_rest 1828.6,
% contrast 1.6133, improvement 17.162 %, interior minimum. The conditioning
% claim is confirmed; the 7.1 % shift in contrast is reported beside it.

\subsection{What the method measures}
\label{sec:discussion:measures}

Section~\ref{sec:constrains} establishes that where this method responds
strongly, the quantity determined is the background density of a body, and
through the identity of Equation~\eqref{eq:excessidentity} the mass anomaly
that accompanies it, rather than the density of the region. It is more
accurate to describe the technique as bounding the mass asymmetry of a body,
from which a density contrast follows only once a dividing surface has been
fixed by other means.

That reframing bears on a question that has stood since 2008, though it does
not dissolve it, and an earlier draft of this subsection claimed more than
the data support. \citet{scheeres2008} inferred a dense neck,
\citet{lowry2014} a dense head, and \citet{kanamaru2019pss} found the neck
model fitting marginally better on the smooth terrains while reporting that
a three-region model is largely unconstrained. The reframing explains why
such a disagreement is possible: a mass anomaly does not specify the region
carrying it, and two structures placing comparable excess mass at comparable
distance from the centre of figure produce comparable centre-of-mass offsets
and comparable dispersion reductions. The degeneracy is the one already
identified in Section~\ref{sec:constrains}: the pair $(\densR,\VR)$ is free
along $\Ddens\VR = \mathrm{const}$ at fixed first moment, and no surface
measurement resolves it.

What must be said, however, is that the head and neck models of
Section~\ref{sec:itokawa:ambiguity} are \emph{not} an instance of that
degeneracy, and citing them as one would misdescribe our own result. They do
not place comparable excess mass: the head model returns
$\Dmass/\Mtot = 7.69\%$ and the neck model $3.47\%$, a factor of $2.2$, and
the slab's centroid lies much closer to the centre of figure, so the two
mass dipoles differ by more still. The objective distinguishes them
decisively, by $16.9\%$ against $2.3\%$ in improvement---a factor of
$7.3$---and by $5.8$ against $2.0$ on the displacement statistic
$\Sstat$. On the global surface the head model is preferred and the neck
model does not clear the resolution envelope at all. The genuine ambiguity
is between structures the method cannot separate; head and neck on Itokawa
are structures it can.

\subsection{An analytic test of the Paper~IV resolution result}
\label{sec:discussion:ellipsoid}

Paper~IV compared the sensitivity of two surface quantities to mesh
resolution, and concluded that the normalized geopotential dispersion is far
the more stable. That comparison, like every comparison of its kind, took the
finest available mesh as the reference. The premise is reasonable but it is
not proof: if both quantities were biased in the same direction, or if the
decimation were introducing a systematic of its own, mesh-to-mesh variation
would not reveal it.

A homogeneous ellipsoid removes the assumption. Its interior gravity, and
hence its surface gravity as the limit from within, has an exact closed form
\citep{macmillan1930}:
%
\begin{equation}
g_i = -2\pi G \varrho\, abc\, A_i x_i , \qquad
A_i = \int_0^\infty \frac{\mathrm{d}u}
     {(a_i^2+u)\sqrt{(a^2+u)(b^2+u)(c^2+u)}} ,
\label{eq:macmillan}
\end{equation}
%
which reduces to $GM/R^{2}$ in the sphere limit. Our evaluation of
Equation~\eqref{eq:macmillan} reproduces that limit to $2.2\times10^{-14}$
in relative terms. That figure is a check on the numerical quadrature of
$A_i$ and on nothing else; it is not a statement about the polyhedral code,
which cannot reproduce a sphere to any such precision, since an inscribed
icosphere of the facet counts used here falls short of $4\pi/3$ in volume by
$0.05\%$ (Section~\ref{sec:methods:verify}). Because the field points of
Section~\ref{sec:methods:model} are displaced inward from the facet
centroids, the interior form of Equation~\eqref{eq:macmillan} is the
appropriate reference at every point compared, and no exterior continuation
is required. The slope at any surface point follows from the analytic gravity
and the analytic normal of the ellipsoid, and the area-weighted means follow
by quadrature on the exact surface. Nothing in the reference depends on a
mesh, on the polyhedral code, or on any decimation.

We tested three ellipsoids, three spin states and five facet counts, ninety
configurations in all. The mean absolute error against the analytic value is
given in Table~\ref{tab:ellipsoid}.

\begin{table}[htbp]
\centering
\small
\caption{Mean absolute error against the analytic ellipsoid, averaged over
three shapes, three spin states and both decimators at each facet count. The
two decimators are separated in Table~\ref{tab:decimator}; the entries here
are their mean, which is why the column means reproduce the decimator
figures quoted below. The final two columns give the error as a fraction of
the effect this paper measures: the dispersion error is between three and
four orders of magnitude smaller than the $16.9\%$ improvement recovered on
Itokawa, which is the figure of merit that matters for the inversion, as
distinct from the accuracy of either quantity in its own right.}
\label{tab:ellipsoid}
\setlength{\tabcolsep}{10pt}
\begin{tabular}{rrrrr}
\toprule
Facets & Mean slope & Dispersion & Ratio
 & dispersion error $/\improve$ \\
\midrule
$2{,}000$  & $11.62\pm3.33\%$ & $0.196\pm0.023\%$ & $59\pm18$ & $1.2\times10^{-2}$ \\
$4{,}000$  & $4.74\pm1.35\%$  & $0.117\pm0.014\%$ & $40\pm13$ & $6.9\times10^{-3}$ \\
$8{,}000$  & $2.30\pm0.63\%$  & $0.049\pm0.004\%$ & $47\pm14$ & $2.9\times10^{-3}$ \\
$16{,}000$ & $0.93\pm0.25\%$  & $0.028\pm0.003\%$ & $33\pm9$  & $1.7\times10^{-3}$ \\
$32{,}000$ & $0.33\pm0.07\%$  & $0.017\pm0.002\%$ & $19\pm4$  & $1.0\times10^{-3}$ \\
\bottomrule
\end{tabular}
\end{table}
% The 49,152-facet row was attempted and abandoned. It needs a source mesh
% finer than the target, and the next subdivision gives 327,680 facets, at
% which a single configuration costs some 256 times the work of one at
% 20,480; nine configurations would have run for about a week. One
% configuration completed before the run was stopped, the triaxial ellipsoid
% without spin: slope errors of +0.15 % (quadric) and +0.19 % (clustering)
% against dispersion errors of -0.02 %, a ratio near 8. It is one shape and
% one spin state and is not comparable with the pooled rows above, so it is
% reported in the text as an indication and not added to the table.

The dispersion is the more accurate at every facet count, by a factor that
runs from $59$ at the coarsest mesh to $19$ at the finest, and the trend
continues beyond the table. At $49{,}152$ facets, the resolution at which the
Itokawa and Eros results are computed, a single configuration gives slope
errors of $+0.15\%$ for quadric collapse and $+0.19\%$ for clustering
against dispersion errors of $-0.02\%$: a ratio near $8$. That is the figure
a summary of this paper should quote, since it is the resolution the results
use, and it carries two qualifications. It is one shape and one spin state
rather than the average of nine, so it has none of the averaging the rows
above have; and the dispersion error it is formed against is quoted to one
significant figure, so the ratio is determined only to within roughly $6$ to
$12$. The direction is what matters:
the advantage of the dispersion narrows as the mesh refines, and the figures
that should be carried into any judgement about the meshes used in this
paper are the smaller ones at the foot of the table, not the $59$ at its
head. Three qualifications belong with those figures.

The advantage narrows with resolution, and the direction is systematic
rather than accidental. Fitting each column against facet count in
logarithms gives $E_{\mathrm{slope}} \propto N^{-1.26}$ and
$E_{\mathrm{disp}} \propto N^{-0.91}$, so the ratio falls as $N^{-0.35}$:
the slope converges faster. Extrapolated to the $49{,}152$-facet meshes on
which the Itokawa and Eros results are computed, the advantage is of order
$16$ to $20$ rather than $59$, and the figure quoted in an abstract should be
the one at the working resolution. The ratio column is noisy, and quantifiably so. With eighteen
configurations per row the standard errors of the mean are large: the ratio
is $59.3\pm18.4$ at $2{,}000$ facets, $46.8\pm13.5$ at $8{,}000$ and
$19.3\pm4.4$ at $32{,}000$. The departure from the monotone trend at
$8{,}000$ facets is under half a standard error and is not a feature. The
spread arises because three ellipsoid shapes and three spin states are
pooled in each row, and the error of a coarse mesh depends strongly on
both. ``$19$ to $59$'' therefore describes the range of the row means, not a
determination of how the ratio varies: only the two endpoints differ by more
than their errors, and the ordering in between is not resolved.
% Computed (n = 18 per row, both decimators pooled):
%   2000  slope 11.62+-3.33  disp 0.1958+-0.0233  ratio 59.3+-18.4
%   4000  slope  4.74+-1.35  disp 0.1175+-0.0139  ratio 40.4+-12.5
%   8000  slope  2.30+-0.63  disp 0.0492+-0.0044  ratio 46.8+-13.5
%  16000  slope  0.93+-0.25  disp 0.0281+-0.0026  ratio 33.2+- 9.4
%  32000  slope  0.33+-0.07  disp 0.0172+-0.0019  ratio 19.3+- 4.4
% The 8,000-facet ratio departs from the trend by less than half a standard
% error and is not a genuine feature.

Paper~IV inferred a factor of roughly $3$ to $14$ from mesh-to-mesh
variation, and the two numbers should not be read as measurements of the
same thing. Mesh-to-mesh variation measures \emph{stability}---how much a
quantity moves when the discretization changes---while
Table~\ref{tab:ellipsoid} measures \emph{accuracy} against a reference
computed independently of the code. A quantity can be stable and biased, or
unstable and unbiased. What the ellipsoid establishes is that Paper~IV's
ordering of the two quantities is not an artefact of taking the finest mesh
as truth, which is the assumption it could not test; it does not establish
that Paper~IV's factor was conservative, and an earlier draft that said so
was comparing across the two meanings.

The ellipsoid also bounds what it can be used for. Decimating from
$32{,}000$ to $2{,}000$ facets moves the dispersion by about $0.18$
percentage points on these smooth bodies, against the $0.9$ to $5.9\%$
Paper~IV measured over a comparable range on real ones---a factor of five to
thirty. Surface roughness, which an ellipsoid does not have, is the
difference. The mesh-resolution envelope used throughout this paper is
therefore a property of rough bodies that the analytic test cannot
reproduce, and the test validates the envelope's ordering and not its
magnitude.

The same experiment answers a related question about the choice of
decimator. Papers~III and~IV reduced their meshes by vertex clustering; this
paper uses quadric edge collapse, because clustering does not preserve
topology and cannot reach small facet targets
(Section~\ref{sec:methods:shapes}). Against the analytic reference the two
are equivalent in accuracy: the paired differences over the ninety
configurations are $+0.12\pm0.20\%$ on the slope and
$+0.019\pm0.012\%$ on the dispersion, neither of them significant, so the
apparent lead of clustering on the dispersion is not established
(Table~\ref{tab:decimator}). The numbers reported in Papers~III and~IV are
therefore not in question, and the change of decimator in this paper is made
for topology rather than accuracy.

\begin{table}[htbp]
\centering
\small
\caption{The two decimators against the analytic ellipsoid, averaged over
all ninety configurations. The differences are $3\%$ on the slope and $29\%$
on the dispersion, both in favour of clustering; whether either is
resolved requires the standard errors, which are not yet reported.}
\label{tab:decimator}
\setlength{\tabcolsep}{12pt}
\begin{tabular}{lrr}
\toprule
Decimator & Mean slope error & Dispersion error \\
\midrule
Vertex clustering     & $3.93 \pm 1.22\%$ & $0.072 \pm 0.009\%$ \\
Quadric edge collapse & $4.05 \pm 1.16\%$ & $0.091 \pm 0.016\%$ \\
\midrule
Paired difference     & $+0.12 \pm 0.20\%$ & $+0.019 \pm 0.012\%$ \\
\bottomrule
\end{tabular}
\end{table}
% Computed (n = 45 each): clustering slope 3.925+-1.223, dispersion
% 0.0722+-0.0088; quadric slope 4.049+-1.156, dispersion 0.0909+-0.0155.
% The configurations pair exactly, so the paired differences are the sharper
% test: quadric minus clustering is +0.0187+-0.0115 on the dispersion, 1.6
% standard errors, and +0.124+-0.200 on the slope, 0.6. Neither is
% significant, which is what makes "equivalent" defensible.

One observation from the same comparison is not explained. On the four real
bodies of the Paper~IV resolution study, the mean slope under clustering
rises monotonically with facet count on every one, by $+21.9\%$, $+14.2\%$,
$+5.8\%$ and $+1.8\%$, while under quadric collapse there is no monotonic
trend. On the smooth ellipsoid neither method drifts: both overestimate the
slope on coarse meshes by the same amount and converge smoothly. Surface
roughness is the natural explanation, but no analytic reference exists for a
rough body and we make no claim about which decimation is closer to the
truth there.
% Reconciled: 1.8-25.6 % is what Paper IV published, from its own mesh
% ladder. Recomputing the same four bodies for this paper gives 1.8-21.9 %
% with vertex clustering and 3.1-6.2 % with quadric collapse. The lower ends
% agree because Gaspra sets both; the upper end differs because the ladders
% differ. The text below now quotes the recomputed range and attributes the
% published one to Paper IV.

\subsection{Relation to Paper~IV}
\label{sec:discussion:paper4}

Paper~IV showed that the normalized geopotential dispersion is far less
sensitive to mesh resolution than the area-weighted mean slope. That is a
necessary condition for the present inference but not a sufficient one: the
minimum of a stable function need not itself be stable.
Sections~\ref{sec:itokawa:mesh} and~\ref{sec:nulls:arrokoth} supply the
missing step, and the form in which they supply it should be stated
precisely, because it is weaker than the bare phrase ``no change'' suggests.

Across the three Itokawa meshes that share a dividing plane the recovered
$\denshead$ is unchanged, which---since no sub-grid refinement is
applied---bounds the variation at half a scan step,
$\SI{12.5}{\kilo\gram\per\cubic\metre}$ or $0.45\%$, rather than resolving
it. On the coarsest mesh the recovered density is not the informative
quantity at all: the detector moves the plane by $\SI{9.4}{\metre}$, which
at the local gradient displaces $\denshead$ by about $3.4\%$ and $\densrest$
by $0.14\%$. Resolution sensitivity enters this calculation through the
plane detector rather than through the objective function, and it enters
asymmetrically, mattering for the quantity the method does not constrain and
not for the one it does. On Arrokoth the recovered density is displaced by no
more than one scan step across three meshes. The two claims---that the
objective is stable, and that its minimizer is---are separate, and both now
have evidence, of different strengths.

One correction to Paper~IV should be recorded here. It attributes the
relaxation reasoning to \citet{kanamaru2019}, who in fact cite
\citet{richardson2014} for it in their own introduction. The argument, the
technique of minimizing potential variance, and the quantity Paper~IV calls
the normalized geopotential dispersion are all due to Richardson and
Bowling; Kanamaru and Sasaki extended the technique from a single mean
density to an interior distribution. Paper~IV's reference list contains no
citation to \citet{richardson2014}, and it should have.

\subsection{One strong signal in five bodies}
\label{sec:discussion:rate}

Whether one strong signal in five bodies reflects the rarity of large
interior contrasts or the sensitivity floor of the method is not settled by
this work, and we do not claim to settle it. The four weak-signal cases are
individually explicable. Eros is measured to be nearly uniform
\citep{yeomans2000,scheeres2020}. Bennu is heterogeneous but radially so
\citep{scheeres2020}, in a geometry a plane cut cannot see.
\citet{jorda2016} state that 67P's centre-of-mass position cannot be
explained by different but individually homogeneous lobe densities, which is
precisely the model class fitted here. Arrokoth has no measured mass, and
its answer sits at the threshold of a sign reversal across the plausible
range of assumptions.

What is missing is a calibrated detection limit: the contrast a body of
given size and spin must carry before this method would see it above the
mesh-resolution sensitivity envelope. The synthetic experiment of
Section~\ref{sec:constrains:synthetic} was designed for a different purpose
and does not supply one, since it varies geometry rather than injecting a
contrast of known size.

Establishing that limit is harder than it appears. Planting a contrast in a
body and leaving its shape alone tests nothing, because the objective
function is built from the shape and the trial density and never reads a
field computed beforehand: the inversion would return whatever the shape
alone would have given. The premise asserts that the surface has relaxed
onto an equipotential of the true interior, so a body with a genuine
contrast is not the same shape as a uniform one, and a ground-truth test
must build the shape the premise describes before inverting it.

We attempted this and do not report the outcome, because the construction did
not converge. Relaxing a synthetic surface toward an equipotential of a
planted two-region interior, by displacing vertices in proportion to the
local departure of the geopotential from its area-weighted mean, reduces the
dispersion but does not settle: on a planted contrast of $1.5$ the residual
fell from $0.066$ to $0.030$ over $140$ iterations and then rose again, and
the relaxed body was no longer a contact binary, the plane holding a given
volume fraction having moved by six times its original offset. Recovered
contrasts from a surface that has not relaxed are not a test of the premise,
and their correlation with the residual---$-0.84$ across the seven rungs of the
planted ladder, at contrasts of $1.00$, $1.05$, $1.10$, $1.25$, $1.50$,
$1.75$ and $2.00$---indicates that the numbers measure the failure of the
relaxation rather than the behaviour of the inversion.
% Seven rungs, listed in the text; the correlation is over n = 7.

The mode of failure is itself informative and points to what a working
scheme would have to do. The plane holding a fixed volume fraction moved as
the surface relaxed, so the region and the shape co-evolve: relaxing toward
an equipotential of a \emph{fixed} interior is not a well-posed target once
the interior is defined by a surface that is itself moving. A fixed point
requires solving for the shape and the dividing surface together, which is
the same coupling that Section~\ref{sec:plane} reports from the other
direction---the objective does not determine its own plane. A stable scheme,
displacing along surface normals with a step controlled to reduce the
dispersion monotonically and with the volume distribution constrained, is a
piece of work in its own right. We regard it as the most useful extension of
the present study, and as a paper rather than an appendix.

A second extension needs no new machinery at all. Every body examined here
was chosen because its mass is already known, or, in Arrokoth's case,
because the consequences of not knowing it are instructive. The reverse case
is available: a body with a resolved shape model and a spin state but no
mass, which a spacecraft will visit and weigh. For such a target the
inversion can be run now as a function of the assumed bulk density, giving
a curve of contrast against assumed bulk density, as in
Section~\ref{sec:nulls:arrokoth}, rather than a number, together with
the amplification factor and the saddle prominence that say how sharply the
answer depends on the assumption. When the mass arrives it selects a point
on a curve published beforehand. That is the one form of external test this
method has not had: the comparisons with \citet{lowry2014} and
\citet{kanamaru2019} in Section~\ref{sec:itokawa} were made against
measurements that already existed when we made them.

\subsection{On verification practice}
\label{sec:discussion:verification}

This series has returned repeatedly to the question of how numerical results
in this area should be checked, and this paper adds the sharpest instance.
The cone construction of Section~\ref{sec:methods:cone} passed a
volume-closure test at $2.1\times10^{-14}$ in relative terms while enclosing
a solid wrong by a factor of $1.6$, because the test compared the mesh
against the same tetrahedra that built it. The test is not vacuous---it would
catch an orientation error or a gross failure of the tetrahedral sum---but it
is blind by construction to the one property that was wrong, which is which
solid the mesh encloses. The predicate test that does distinguish the two
solids costs three lines of code.

The other failures of Section~\ref{sec:failures} share a property: each was
invisible in the configuration that mattered and appeared only when a
parameter was varied that the analysis did not require. The NaN failure was
total at one plane position and absent at another $\SI{30}{\metre}$ away.
The edge minima appeared only when a scan was widened. The misplaced
dividing plane was found by comparing two selection rules on the same
cross-sectional profile---the saddle criterion of
Equation~\eqref{eq:prominence} against the absolute interior minimum---and
noticing that they disagreed by $\SI{15}{\kilo\metre}$; a study committed to
either rule alone would not have discovered it. A study reporting one run at
one plane on one mesh would have met none of these and would have reported
plausible numbers throughout. In this method, parameter sweeps are not
supplementary material; they are the only way the failure modes become
visible.

The buried facets of Section~\ref{sec:failures:buried} make a sharper
version of the same point. The manifold criterion was added precisely to
catch what a volume comparison cannot, and it passes them without complaint,
because a mesh in two closed pieces is topologically sound whatever their
spatial arrangement. Each new check inspects the property its author had in
mind and no other, and on this occasion the omission would have carried a
qualitative conclusion: a sign reversal that is not there.

The analytic ellipsoid of Section~\ref{sec:discussion:ellipsoid} illustrates
the complementary point. A reference computed independently of the code under
test can settle in one experiment what mesh-to-mesh comparison can only
suggest, and it can overturn a conclusion drawn from such comparison: the
provisional reading of the decimator evidence, that quadric collapse is the
more accurate, did not survive contact with it.

\subsection{Limitations}
\label{sec:discussion:limits}

The model admits two uniform regions divided by a plane perpendicular to a
principal axis. Real interiors need not take that form, and for two of the
bodies here it is known that they do not: \citet{jorda2016} exclude it for
67P, and \citet{scheeres2020} report a radial structure on Bennu that the
parameterisation cannot represent. \citet{kanamaru2019pss} show that
relaxing the model to three regions on Itokawa leaves it largely
unconstrained, so enrichment is not a straightforward remedy.

The relaxation premise itself is unverified on every body examined, and
fails its first stated condition on the one that yields a strong signal
(Section~\ref{sec:discussion:premise}). Section~\ref{sec:discussion:conditioning}
argues that this condition governs the conditioning of a different inverse
problem and does not transfer to a fixed-mass two-region inversion, but that
argument is structural: it removes a reason to expect failure and does not
establish that the surface has relaxed.

The dividing plane rests on judgement informed by two external criteria
rather than on anything measured within the method
(Section~\ref{sec:plane:external}), and the numerical diagnostics respond to
a misplaced one unevenly (Section~\ref{sec:failures:plane}). Of the four
tested, the amplification factor fails outright, ranking the misplaced plane
as the strongest signal in the study; the displacement statistic ranks it
correctly at $\Sstat = 2.86$ against $5.8$ but places it above all four
genuine nulls; the envelope comparison rejects it by a margin of $6\%$, and
only if the body-specific envelope value is above the $5.526\%$ improvement
it returns; and the decisive argument is the provenance of the plane, which
is not available from the output.

The recovered densities are quantized at the $\SI{25}{\kilo\gram\per
\cubic\metre}$ scan step, with no sub-grid refinement
(Section~\ref{sec:methods:minimisation}). The consequences are visible
throughout---two planes $\SI{0.8}{\metre}$ apart returning the same
$\denshead$, two non-monotone entries in the $\dCOM$ column, and an
improvement that can come out marginally negative where the grid does not
contain the uniform solution---and one of them, reported in
Section~\ref{sec:failures:small}, locates the true minimizer at a plane where
the grid reports a value $\SI{12}{\kilo\gram\per\cubic\metre}$ away. Solving
the stationarity condition of Equation~\eqref{eq:affine} in closed form would
remove the grid from the error budget entirely, since $\potvarn^{2}$ is a
ratio of quadratics in $\densR$ and its stationarity condition is linear with
a single root. We have not done so, and it is the clearest improvement
available to a future implementation.

The synthetic calibration of the background-density invariance rests on four
bodies carrying a measurable exponent, three of them clustered in signal
strength, so the direction of the effect is supported while its functional
form is not determined; and because those bodies are uniform, the
identification of a strong inversion response with a physical mass anomaly
remains an interpretation (Section~\ref{sec:constrains:limits}).

One shape model required a correction that others did not. The Arrokoth
model is supplied as two unjoined shells with $0.943\%$ of the surface area
lying inside the opposite shell, and excluding those facets changes the
recovered contrast at the detected plane, under the head model, from
$1.062$ to $0.9999$ at the lowest assumed bulk density
(Section~\ref{sec:nulls:arrokoth}). No other model in the study required
this, but it is not a defect peculiar to Arrokoth: any shape model delivered
as more than one component may have the same property, and the test costs one
ray-casting pass.

Finally, we have not obtained \citet{motooka2012}, which
\citet{kanamaru2019} describe as an earlier attempt to fit the global shape
of Itokawa to a single equipotential surface, and which may therefore be
closer to the approach taken here than any other prior work. It is cited as
described by Kanamaru and Sasaki, and no claim is made about its contents.

Section~\ref{sec:guidelines} sets out what these limitations imply for
anyone applying the method.
